% =====================================================================
%  IRV : 오류주입공격과 감염형 무작위 검증 대응기법
%  Overleaf 업로드용 마스터 파일 (한글 초안, kotex).
%  구조:  main.tex  +  sections/*.tex  +  figures/*  +  references.bib
%  컴파일: pdfLaTeX (kotex) → BibTeX → pdfLaTeX ×2
%  그림은 figures/make_figures.py 로 생성: fig_coverage, fig_cost (2차오류는 표로 대체)
% =====================================================================
\documentclass[11pt]{article}

\usepackage[english]{babel}
\usepackage{kotex}                 % 한글
\usepackage{amsmath,amssymb,bm}
\usepackage{graphicx}
\usepackage{booktabs}              % \toprule \midrule \bottomrule
\usepackage{float}                 % [H] : 표를 정의 위치에 고정(챕터 밀림 방지)
\usepackage{algorithm}
\usepackage{algpseudocode}
\usepackage{textcomp}
\usepackage[letterpaper,top=2.3cm,bottom=2.3cm,left=2.7cm,right=2.7cm]{geometry}
\usepackage[colorlinks=true, allcolors=blue]{hyperref}

\graphicspath{{figures/}}

\newcommand{\zvec}{\mathbf{z}}
\newcommand{\yvec}{\mathbf{y}}
\newcommand{\svec}{\mathbf{s}}
\newcommand{\Amat}{\mathbf{A}}

\title{격자 기반 전자서명 HAETAE에 대한\\
       오류주입공격과 감염형 무작위 검증 대응기법}
\author{Philip Johansson}
\date{2026}

\begin{document}
\maketitle

\begin{abstract}
양자 위협에 대응하여 제안된 격자 기반 Fiat--Shamir 서명 스킴 \texttt{HAETAE}는 수학적 안전성과
별개로 실제 구현 환경에서 오류주입공격(Fault Injection Attack, FIA)에 취약할 수 있다. 특히 서명의
응답 연산 $\zvec=\yvec+(-1)^b c\svec_1$은 비밀키와 일회용 난수가 결합되는 유일한 지점으로,
이 연산의 무결성이 보장되지 않으면 비밀 정보가 직접 노출될 수 있다. 본 논문에서는 \texttt{HAETAE}의
서명 구조를 분석하여 응답 연산에 대한 세 가지 오류 지점($c\svec$ 스킵, $+\yvec$ 스킵, 거부 우회)을
도출하고, 그 중 $+\yvec$ 덧셈을 건너뛰는 경우 단 한 번의 오류만으로 비밀키 $\svec_1$을 직접
복원할 수 있음을 보인다. 전체 서명 구현에서
(ARM 기반 MCU에 에뮬레이션된 단일 오류 모델 하에) 오류가 주입된 단일 서명으로부터 $\svec_1$의 768개 계수 전부를
100\% 복원함을 확인하였으며, 복구 파이프라인의 정확성은 자체 검증으로 입증하였다. 또한 본
논문은 연산 재계산 검증, 서명 경계 노름 재검사, 비밀키 무결성 검사를
비교(분기) 연산에 의존하지 않는 감염형 마스킹으로 통합한 경량 대응 기법
(IRV, Infective Randomized Verification, 감염형 무작위 검증)를
제안한다. 제안 기법은 (i)~거부 우회를 포함한 7개 오류 지점을 모두 차단하고, (ii)~탐지 분기를
건너뛰는 2차 오류에 내성을 가지며, (iii)~전체 이중 연산 대비 절반 수준의 서명 시간으로
동작함을 보인다.

\end{abstract}

% 영문 초록 (국문 학술지 제출 규정 대응; 필요 없으면 아래 블록 삭제)
{\renewcommand{\abstractname}{Abstract}%
\begin{abstract}
\texttt{HAETAE}, a lattice-based Fiat--Shamir signature scheme proposed for the
post-quantum era, may be vulnerable to fault injection attacks (FIA) in real
implementations regardless of its mathematical security. In particular, the response
computation $\zvec=\yvec+(-1)^b c\svec_1$ is the only step in which the secret key and
the one-time nonce are combined; if its integrity is not guaranteed, secret information
can be exposed directly. We analyze the \texttt{HAETAE} signing procedure and identify
three fault points in the response computation (skipping $c\svec$, skipping $+\yvec$, and
rejection bypass). We show that skipping the $+\yvec$ addition enables the secret key
$\svec_1$ to be recovered directly from a \emph{single} faulted signature. On a
full-signature STM32F405 implementation with ChipWhisperer-Husky, under an
\textbf{emulated single-fault model}, we recover all 768 coefficients of $\svec_1$ with
100\% success from one faulted signature, and the recovery pipeline is confirmed by a
self-test. We further propose \texttt{IRV}
(Infective Randomized Verification), a
lightweight countermeasure that unifies recomputation-based verification, signing-boundary
norm re-checking, and secret-key integrity checking into \textbf{branchless infective
masking} that does not rely on comparison (branch) operations. Experimentally, IRV
(i)~blocks all seven fault points including rejection bypass, (ii)~resists second-order
faults that skip detection branches, and (iii)~operates at less than half the
\emph{signing time} of full double computation ($+14\%$ signing-time overhead), though it
requires more code/RAM than double computation.

\smallskip
\noindent\textbf{Keywords.} Post-quantum cryptography; Lattice-based signature;
\texttt{HAETAE}; Fault injection attack; Infective countermeasure; Single-signature key recovery.

\end{abstract}}

\section{서론}
\label{sec:intro}

양자컴퓨팅 기술의 발전은 소인수분해 및 이산로그 문제에 기반한 기존 공개키 암호체계의 안전성을
근본적으로 위협하고 있다. 이에 따라 미국 국립표준기술연구소(NIST)는 2024년 양자내성암호
(Post-Quantum Cryptography, PQC) 디지털 서명 표준으로 \texttt{CRYSTALS-Dilithium}(ML-DSA)을
선정하였으며~\cite{dilithium,fips204,nist_pqc}, 국내에서도 KpqC 공모를 통해 격자 기반 서명
\texttt{HAETAE}가 최종 후보로 선정되었다. \texttt{HAETAE}~\cite{haetae}는 Module-LWE 및
Module-SIS 문제에 안전성을 기반하는 Fiat--Shamir with Aborts 방식의 격자 서명이다. Gaussian
샘플링에 의존하던 기존 스킴과 달리, 서명 응답 벡터를 고차원 $\ell_2$-노름 공(\emph{hyperball})
상의 균등분포에 맞추어 거부 샘플링(rejection sampling)하고, 여기에 \texttt{BLISS}에서 도입된
\emph{bimodal} 기법(부호 $b\in\{0,1\}$로 응답의 방향을 뒤집어 거부율을 낮추는 기법)을 결합한다.
이 hyperball 균등 샘플링을 고정소수점(fixed-point) 산술로 구현함으로써, 서명 및 공개키 크기를
효과적으로 줄이면서도 구현 효율을 확보한 것이 특징이다.

그러나 암호 알고리즘의 수학적 안전성은 실제 하드웨어 구현 환경에서 그대로 보장되지 않는다.
성능 최적화와 구현 제약으로 인해 새로운 공격 표면이 발생할 수 있으며, 이는 클럭
글리치~\cite{yang2020}, 전압 글리치~\cite{oflynn2016}, 레이저~\cite{skorobogatov2002}, 전자기파
~\cite{dumont2019} 등을 이용한 오류주입공격(Fault Injection Attack, FIA)에 취약해질 수
있다~\cite{boneh1997,barel2006,espitau2016}. 오류주입공격은 암호 연산의 흐름이나 내부 데이터를
왜곡하여 비밀 정보를 추출하는 강력한 능동적 공격이며, 단 한 번의 오류만으로도 전체 시스템의 안전성이
위협받을 수 있으므로 구현자는 이를 심각하게 고려하여야 한다~\cite{ravi2022survey}. 이러한 위협은
RSA~\cite{lenstra1996}, 타원곡선 암호~\cite{biehl2000,ciet2005} 등 기존 공개키 암호 체계 전반에서
이미 확인된 바 있다. 실제로 \texttt{Dilithium}에 대해서도 단일 또는 소수의 오류로 비밀 키의 전부
또는 일부를 복구할 수 있음이 실험적으로 입증되었다~\cite{ravi2019determinism,krahmer2024,elghamrawy2023}.

\texttt{HAETAE}의 서명에서 가장 핵심적인 연산은 챌린지 $c$에 대한 응답 벡터를 계산하는
$\zvec=\yvec+(-1)^b c\svec_1$이다. 여기서 $\yvec$는 일회용 난수, $b$는 bimodal 부호, $\svec_1$은
비밀키 벡터로, 이 식은 서명 과정에서 비밀키를 직접 사용하는 유일한 단계이므로 오류주입의 주된
표적이 된다. 그럼에도 \texttt{HAETAE}의 전체 서명 과정을 대상으로 한 응답 연산 오류 분석과 이에
대한 경량 대응 연구는 아직 충분히 보고되지 않았다.

이에 본 논문에서는 \texttt{HAETAE}의 결정론적 서명 구조를 분석하여 응답 연산에 대한 오류주입공격을
체계적으로 규명하고, 이를 효율적으로 방어하는 대응 기법을 제안한다. 본 논문의 기여는 다음과 같다.
\begin{itemize}
  \item \texttt{HAETAE} 응답 연산에 대한 오류 분석을 제시하고, 결정론 격자 서명의 skip-addition
        계열 공격~\cite{ravi2019determinism}을 HAETAE 응답에 특화하여 \textbf{$+\yvec$ 스킵이 단일
        서명으로 $\svec_1$을 직접 복원}시킴을 수학적으로 보인다. 이를 STM32F405 전체 서명에서
        소프트웨어 오류 주입(\textbf{에뮬레이션 단일 오류 모델}) 하에 실증한다($\svec_1$ 768계수 100\%
        복원(가역 챌린지 조건), 복구 파이프라인은 self-test로 검증).
  \item 기존 격자 오류 대응이 모두 \emph{탐지 후 중단(detect-and-abort)} 방식인 것과 달리,
        \textbf{비교(분기) 연산에 의존하지 않는 감염형 무결성 보호 원리}를 격자 서명에 확장한
        대응 기법 \textbf{IRV}를 제안한다(Seo 등의 RSA-Fermat 감염형 대응~\cite{seo2013}을 발전시킨 것).
        이는 \emph{우리가 아는 한 격자 Fiat--Shamir 서명에 대한 최초의 무분기 감염형 대응}이다.
  \item 단순한 서명 후 검증의 한계를 규명한다. \texttt{HAETAE}의 검증 경계 $B_2$는 서명 경계
        $B_1/B_0$보다 느슨하여($B_1<B_2$) 거부 우회 오류를 통과시키며, IRV의 서명 경계 재검사가
        이 사각지대를 차단한다.
  \item 탐지 분기를 건너뛰는 2차 오류(탐지 분기 스킵 모델)에 대한 내성을 실험적으로 입증한다. 이중
        연산과 서명 후 검증은 단일 분기 스킵으로 우회되나, 비교 연산이 없는 IRV는 무력화되지 않는다.
  \item 소프트웨어 라인 주입(대응 기법의 공정 비교)을 통해 커버리지,
        누설 방출률, 오버헤드, 연산 시간을 정량적으로 비교한다.
\end{itemize}

본 논문의 구성은 다음과 같다. \ref{sec:bg}장에서는 \texttt{HAETAE}와 오류주입공격의 배경을
설명하고, \ref{sec:fault}장에서는 응답 연산의 오류 지점 분석과 단일 서명 비밀키 복원 공격을 다룬다.
\ref{sec:irv}장에서는 제안 기법 IRV를 제시하고, \ref{sec:eval}장에서 실험적으로 평가한 뒤 한계를
논의한다. \ref{sec:related}장에서 관련 연구를 다루고 \ref{sec:concl}장에서 결론을 맺는다.

\section{배경}
\label{sec:bg}

\subsection{HAETAE 서명과 Fiat--Shamir with Aborts}
Shor 알고리즘~\cite{shor1994}에 의한 양자 위협으로 격자 기반 암호가 차세대 공개키 암호의 유력한
대안으로 부상하였으며~\cite{peikert2016}, 그 중 trapdoor 없는 Fiat--Shamir 격자
서명~\cite{lyubashevsky2012}이 표준화의 중심에 있다. Fiat--Shamir with
aborts~\cite{lyubashevsky2009}는 비밀키 $\svec$, 챌린지 $c$, 일회용 난수
$\yvec$로부터 응답 $\zvec=\yvec+c\svec$를 계산하고, $\zvec$가 비밀 정보를 노출할 위험이 있으면
해당 시도를 거부(abort)하고 새로운 난수로 서명을 재시도하는 거부 샘플링 구조를 사용한다.
\texttt{HAETAE}~\cite{haetae}는 여기에 bimodal 부호 $b\in\{0,1\}$로 $\pm c\svec$를 선택하고
초구 균등 분포에서 난수를 샘플링하여 서명 크기를 줄인다~\cite{bliss2013}. 다항식 연산은
$R_q=\mathbb{Z}_q[x]/(x^{256}+1)$, $q=64513$ 위에서 NTT로 수행되며, 본 논문은 보안 수준
120(MODE2: $L=4$, 비밀 다항식 $M=3$)을 기준으로 한다. 이때 거부 샘플링은 서명 경계 $B_1$(및
보조 bimodal 경계 $B_0$)을 사용하는 반면, 검증은 별도의 경계 $B_2$를 사용하며 \textbf{$B_0,B_1<B_2$}가
성립한다. 이 경계 차이는 \ref{sec:irv}장에서 단순 서명 후 검증의 한계를 설명하는 핵심 근거가 된다.

\subsection{오류주입공격 모델}
오류주입공격은 클럭/전압 글리치, 전자기파, 레이저 등을 통해 연산 중 디바이스의 상태를 일시적으로
교란하여, 출력된 오류 서명을 분석함으로써 비밀 정보를 복원하는 능동적 물리 공격이다. 오류 공격자 모델의 정형화는 이론과 실제를 잇는 관점에서 활발히 논의되어
왔다~\cite{toprakhisar2024,spruyt2021}. 본 논문은 \textbf{단일 오류} 모델을 가정하며, 그 효과를
(i)~명령어 스킵, (ii)~루프 조기 종료, (iii)~데이터 변조의 세 가지로 모델링한다. 이는 오류주입분석에서 널리 채택되는 표준 모델이다~\cite{espitau2016,bindel2016,ravi2022survey}. 공격자는 오류 상태에서 방출된 서명만
관찰할 수 있으며, 비밀키와 난수는 직접 관찰하지 못하고 챌린지 $c$ 등 공개 정보만을 이용한다. 일부
분석에서는 대응 기법의 검사 분기를 무력화하기 위한 \textbf{2차 오류}까지 고려한다.

\subsection{감염형 대응 기법의 종류}
오류 검출의 통상적 방어는 출력 직전에 결과를 검사하는 것이나, Yen과 Joye~\cite{yen2000}는 ``출력
전 검사''만으로는 오류 기반 분석을 충분히 막지 못함을 보였다. 감염형(infective) 대응 기법은 오류를
탐지한 뒤 연산을 중단하는 대신, 오류의 잔차를 출력에 확산시켜 결과 자체를 무효화한다. 즉 ``오류가
발생하면 중단한다''는 \emph{비교(분기) 연산}을 두지 않으므로, 그 분기를 건너뛰는 2차 오류에
본질적으로 강하다. Seo 등~\cite{seo2013}은
Fermat 소정리를 이용하여 비교 연산 없이 적은 추가 연산만으로 RSA 오류를 검출하는 기법을
제안하였으며, 본 논문은 이 설계 철학을 격자 기반 Fiat--Shamir 서명으로 확장한다.

\section{HAETAE 응답 연산 오류 분석과 공격 실현}
\label{sec:fault}\label{sec:attack}

\subsection{공격 대상 및 위협 모델}
본 논문이 주목하는 공격 대상은 거부 샘플링을 통과하여 실제로 방출되는 응답 벡터
$\zvec=\yvec+(-1)^b c\svec_1$의 계산 과정이다. 공격자는 이 연산이 수행되는 시점에 단일 오류를
주입하고, 그 결과 출력된 서명을 관찰한다. 공격자는 방출된 서명(챌린지 $c$ 포함)만 관찰할 수 있으며,
비밀키와 일회용 난수는 직접 관찰하지 못한다.

\subsection{HAETAE 서명 알고리즘}
\label{subsec:sign-alg}
오류 지점을 정확히 규정하기 위해 먼저 \texttt{HAETAE}의 서명 절차를 정리한다(알고리즘~\ref{alg:sign}).
서명은 거부 샘플링 루프로, 매 시도에서 (1)~초구(hyperball)에서 일회용 난수 $\yvec=(\yvec_1,\yvec_2)$와
bimodal 부호 $b$를 샘플링하고, (2)~$\lfloor\yvec\rceil$로부터 챌린지 $c$를 유도하며(상위비트와
$\textsc{LSB}$로 결정), (3)~응답 $\zvec=\yvec+(-1)^b c\svec_1$을 계산하고, (4)~서명 경계 $B_1$과
bimodal 경계 $B_0$에 대한 거부 판정을 수행한다. 판정을 통과하면 서명 $\sigma=(c,\zvec,\mathbf{h})$를
방출하고, 실패하면 새 난수로 재시도한다. 여기서 3단계는 서명 전 과정에서 \textbf{비밀키 $\svec_1$이
사용되는 유일한 단계}이며, 따라서 오류주입의 핵심 표적이 된다.

\begin{algorithm}[H]
\caption{\texttt{HAETAE}.Sign (MODE2: $L{=}4$, $K{=}2$, $q{=}64513$, $L_n{=}2^{13}$). 오른쪽 주석은 오류 지점.}
\label{alg:sign}
\begin{algorithmic}[1]
\State $(\Amat,\svec_1,\svec_2) \gets \textsc{UnpackSK}(sk)$
      \Comment{\textbf{UNPACK}: $\Amat$ 변조}
\State $seed \gets \textsc{Squeeze}(key,\mu,rnd)$
      \Comment{\textbf{SEED}: $seed{\gets}0$ 고정}
\Repeat
  \State $(\yvec_1,\yvec_2,b) \gets \textsc{SampleHyperball}(seed,\mathit{ctr})$
        \Comment{\textbf{SIGNBIT}: $b$ 반전}
  \State $c \gets \textsc{Challenge}\big(\textsc{HighBits}(\Amat\lfloor\yvec_1\rceil{+}2\yvec_2),\,\textsc{LSB}(\lfloor\yvec_1\rceil_0),\,\mu\big)$
        \Comment{\textbf{LSB}: 다른 $c$}
  \State $(c\svec_1,c\svec_2) \gets (c\cdot\svec_1,\; c\cdot\svec_2)$ \quad(NTT 성분별 곱)
        \Comment{\textbf{T1}: 스킵 $\Rightarrow \zvec{\approx}\yvec$}
  \State $(c\svec_1,c\svec_2) \gets (-1)^b\,(c\svec_1,c\svec_2)$
  \State $(\zvec_1,\zvec_2) \gets (\yvec_1{+}c\svec_1,\; \yvec_2{+}c\svec_2)$ \quad(고정소수점)
        \Comment{\textbf{T2}: $+\yvec$ 스킵 $\Rightarrow \zvec{=}(-1)^b c\svec_1$}
  \State $\mathit{reject} \gets \big(\lVert\zvec\rVert^2 > B_1^2 L_n^2\big) \,\lor\, \neg\,\textsc{Bimodal}_{B_0}(\zvec,\yvec,b)$
        \Comment{\textbf{RB}: 판정 스킵 $\Rightarrow$ 경계초과 $\zvec$}
\Until{$\neg\,\mathit{reject}$}
\State \Return $\sigma \gets \textsc{Encode}(c,\zvec,\mathbf{h})$
\end{algorithmic}
\end{algorithm}

\subsection{오류 지점과 정보 누설}
\label{subsec:faultpoints}
알고리즘~\ref{alg:sign}의 각 연산은 잠재적 오류 표적이다. 서명 경로를 분석하여 7개 오류 지점을
도출하였으며(요약은 표~\ref{tab:faultpoints}), 각 지점의 오류 효과와 그로 인한 누설을 함께 정리한다.
이 중 네 가지는 응답 \emph{이전} 단계를 겨냥하며 선행 연구~\cite{lee2026haetae}에서 제시된 지점으로,
대체로 정상 서명과의 \emph{차분}(둘 이상의 서명)이나 예측 가능해진 난수를 통해 비밀 정보에 접근한다.
\begin{itemize}
  \item \textbf{UNPACK}(1행): 공개 행렬 $\Amat$의 언패킹을 변조 $\to$ $\svec_1$ 차분 누설.
  \item \textbf{SEED}(2행): 샘플링 시드를 0으로 고정 $\to$ $\yvec$가 예측 가능해져 난수 노출.
  \item \textbf{SIGNBIT}(4행): bimodal 부호 비트 $b$ 반전 $\to$ $\svec_1$ 차분 누설.
  \item \textbf{LSB}(5행): $\textsc{LSB}$ 추출을 건너뛰어 다른 챌린지 $c$ 유도 $\to$ $\svec_1$ 차분 누설.
\end{itemize}
나머지 세 가지(T1, T2, RB)는 \textbf{응답 연산 자체}를 직접 겨냥하는 본 논문의 지점으로, 더 강한 누설을
만든다.
\begin{itemize}
  \item \textbf{T1 (CS)}(6행): $c\svec$ 곱셈을 건너뛰면 $\zvec\approx\yvec$가 되어 일회용 난수가 노출된다.
        T2가 난수를 제거하는 것과 \emph{정반대}로 마스킹을 깨는 쌍이며, 노출된 난수를 이용한 복원(차분)은
        본 논문의 범위를 벗어나 후속 연구로 남긴다.
  \item \textbf{T2 (ADDY, 가장 치명적)}(8행): $+\yvec$ 덧셈을 건너뛰면 난수 마스크가 제거되어
        $\zvec=(-1)^b c\svec_1$만 남는다. 챌린지 $c$가 공개이므로 NTT 역산으로 $\svec_1$을 \textbf{단일
        서명에서 직접} 복원한다(\ref{subsec:key-recovery}절).
  \item \textbf{RB (REJECT)}(9행): 거부 판정을 건너뛰면 서명 경계 $B_1/B_0$를 위반한 $\zvec$가 방출된다.
        검증 경계 $B_2$가 더 느슨하여($B_0,B_1<B_2$) 단순 서명 후 검증은 이를 통과시키므로
        (\ref{subsec:b1b2}절), 거부 샘플링 제거로 전사(transcript) 분포가 $\svec_1$에 의존하는
        \emph{통계적} 누설이 된다.
\end{itemize}
응답 연산에서 덧셈이나 곱셈을 건너뛰는 skip-addition 오류는 결정론 격자 서명에 대해 이미 그 유효성이
보고된 바 있으며~\cite{ravi2019determinism}, 본 논문이 주목하는 가장 치명적인 시나리오는 $+\yvec$ 덧셈이
누락되는 T2이다.

\subsection{오류 지점의 소프트웨어 구현}
\label{subsec:fault-impl}
대응 기법을 \emph{공정하게 비교}하려면 확률적 하드웨어 글리치로는 보장할 수 없는 ``모든 변형에 동일한
오류''라는 조건이 필요하다. 이를 위해 각 오류 지점을 \texttt{-DFAULT\_SIM} 빌드에서 슈도코드 라인 단위로
결정론적으로 주입한다. 호스트가 결함 라인을 지정하면(\texttt{'f'} 커맨드), 서명 1회에 대해 해당 라인에서
\emph{단일 일시(one-shot transient)} 오류가 발화한다. 무방어(\textsf{baseline})와 세 대응
변형(\textsf{double}/\textsf{leeha}/\textsf{IRV})이 \emph{동일한 훅}을 공유하므로, 같은 오류가 같은 위치에
주입되어 \ref{sec:eval}장의 커버리지 비교가 성립한다. 각 지점의 실현은 다음과 같다(표~\ref{tab:faultpoints}의
``구현'' 열).
\begin{itemize}
  \item 응답 이전 지점은 대상 값을 무력화한다: 시드 버퍼를 0으로(\textbf{SEED}), 언패킹된 $\Amat$를
        0으로(\textbf{UNPACK}), $\textsc{LSB}$ 결과를 0으로(\textbf{LSB}), 부호 비트를 반전(\textbf{SIGNBIT}).
  \item \textbf{T1}: 곱 결과 $c\svec_1,c\svec_2$를 0으로 덮어써 곱셈 스킵을 재현한다($\zvec\approx\yvec$).
  \item \textbf{T2}: $+\yvec$ 덧셈에서 난수 항 $\yvec$를 0으로 대체한다($\zvec=(-1)^b c\svec_1$).
  \item \textbf{RB}: 거부 분기(\texttt{if reject goto \dots})를 실행하지 않아 경계 초과 $\zvec$를 방출한다.
\end{itemize}
또한 대응 기법 평가를 위해, 탐지-후-중단(detect-and-abort) 분기를 겨냥하는 \textbf{2차 오류} 훅을 별도로
두어 대응 기법의 검사 분기 스킵을 모델링한다(\ref{sec:eval}장). 이 결정론적 주입은 단일 일시 오류의
\emph{논리적} 효과를 분리하며, 그 효과는 32비트 마이크로컨트롤러에서 물리적으로 실증된 명령어 스킵
모델과 부합한다~\cite{moro2013}.

\begin{table}[H]\centering
\caption{\texttt{HAETAE} 서명의 오류 지점. ``구현''은 축약형 \texttt{-DFAULT\_SIM} 실현이며 $\dagger$ 표시는
본 논문에서 도출한 지점이다.}
\label{tab:faultpoints}
\footnotesize
\begin{tabular}{@{}lllll@{}}
\toprule
ID & 오류 효과 & 구현 & 누설 정보 & 출처 \\
\midrule
SEED                 & 시드 0 고정 $\to$ 예측 가능한 $\yvec$ & $seed{\gets}0$ & 난수 & \cite{lee2026haetae} \\
SIGNBIT              & 부호 비트 $b$ 반전 & $b{\gets}b\oplus1$ & $\svec_1$(차분) & \cite{lee2026haetae} \\
UNPACK               & 공개 행렬 $\Amat$ 변조 & $\Amat{\gets}0$ & $\svec_1$(차분) & \cite{lee2026haetae} \\
LSB                  & \texttt{poly\_lsb} 스킵 $\to$ 다른 $c$ & $\textsc{lsb}{\gets}0$ & $\svec_1$(차분) & \cite{lee2026haetae} \\
T1 (CS)$^\dagger$    & $c\svec$ 곱셈 스킵 $\to \zvec\approx\yvec$ & $c\svec{\gets}0$ & 난수 & 본 논문 \\
T2 (ADDY)$^\dagger$  & $+\yvec$ 스킵 $\to \zvec=(-1)^b c\svec_1$ & $\yvec{\gets}0$ & \textbf{$\svec_1$ 직접} & 본 논문 \\
RB (REJECT)$^\dagger$& 거부 판정 스킵 $\to$ 경계 초과 $\zvec$ & 분기 미실행 & 통계적 키 정보 & 본 논문 \\
\bottomrule
\end{tabular}
\end{table}

\subsection{실험 환경}
\label{subsec:setup}
실험은 ChipWhisperer-Husky~\cite{oflynn2014}와 CW308 보드에 장착한 \textbf{STM32F405}(Cortex-M4,
192\,KB SRAM) 타겟에서 수행하였다. 재현성을 위해 고정 시드 결정론
난수 발생기를 사용하여 동일한 키와 난수로 서명을 반복하였다(배포용이 아닌 실험용 설정이다). 전체
HAETAE-120 서명을 단일 펌웨어로 타겟에 탑재하였다.

오류는 \texttt{-DFAULT\_SIM} 소프트웨어 라인 주입으로 주입한다. 이는 동일한 오류를 동일한 위치에
결정론적으로 주입하여 대응 기법을 \emph{공정하게 비교}하기 위함이며, 단일 일시 오류의 \emph{논리적}
효과를 분리하여 ``모든 변형에 동일한 오류''라는 조건을 성립시킨다.

\subsection{T2 오류를 통한 단일 서명 비밀키 복원}
\label{subsec:key-recovery}

\paragraph{복원 직관(신호 관점).} 서명 응답 $\zvec=\yvec+(-1)^b c\svec_1$은 \emph{비밀 신호} $c\svec_1$이
\emph{일회용 잡음} $\yvec$에 가려진 형태로 볼 수 있다. 정상 서명은 매번 다른 $\yvec$가 신호를 가려
$\svec_1$을 드러내지 않으나, T2 오류는 $+\yvec$를 건너뛰어 \emph{잡음을 통째로 제거}해 $\zvec=\pm c\svec_1$만
남긴다. 다항식 곱 $c\svec_1$은 순환 합성곱이므로 합성곱 정리에 따라 NTT(유한체 위의 FFT) 영역에서
\emph{성분별 곱}이 되고, 따라서 \emph{공개} 챌린지의 스펙트럼 $\hat c$로 나누는(디컨볼루션) 것만으로
$\svec_1$이 복원된다. 게다가 고정소수점 산술($L_n=2^{13}$) 덕분에 T2가 만드는 값은 $L_n$의 정확한
배수여서, 서명에서 $c\svec_1$이 양자화 오차 없이 그대로 읽힌다.

구체적으로 공격자는 다음 단계로 단일 오류 서명에서 $\svec_1$을 복원한다.
\begin{description}
  \item[단계 1)] 응답 계산 시점에 $+\yvec$ 덧셈을 건너뛰도록 오류를 주입한다. 내부 고정소수점 값은
        $L_n(c\svec_1)$이나, 서명에 \emph{방출}되는 값은 반올림 $\tilde\zvec=\lfloor z_1+L_n/2\rfloor\gg13
        = c\svec_1\bmod q$이다($L_n=2^{13}$). T2에서 $z_1$은 $L_n$의 정확한 배수이므로 이 반올림은
        무손실이며, 공격자는 별도의 나눗셈 없이 $c\svec_1\bmod q$를 직접 관측한다.
  \item[단계 2)] 챌린지 $c$는 서명에 함께 방출되어 공개이다. NTT가 합성곱을 성분별 곱으로 보내므로,
        $\hat c=\mathrm{NTT}(c)$의 성분이 0이 아니면 각 비밀 다항식 $\svec_1^{(j)}$($j=1,2,3$)을
        다음과 같이 복원한다.
        \begin{equation}
          \hat\svec_1^{(j)}[k]=\mathrm{NTT}(\tilde\zvec^{(j)})[k]\cdot \hat c[k]^{-1}\pmod q,
          \label{eq:t2recover}
        \end{equation}
        여기서 $\hat c[k]^{-1}$는 $\bmod\,q$ 곱셈 역원이다. $q=64513$이 소수라도 $\hat c[k]=0$인 성분이
        존재할 수 있으며($c$는 비영 계수가 $\pm1$인 고정 무게 다항식), 그 확률은 서명당 대략 $256/q\approx0.4\%$이다.
        해당 성분에서는 정보가 소멸하나, $\svec_1$의 계수가 $\{-1,0,1\}$의 세 값만 가지므로($\eta=1$) 나머지
        성분의 평가값으로 계수영역에서 유일하게 결정하거나 두 번째 서명으로 보완할 수 있다.
  \item[단계 3)] 부호 $(-1)^b$와 ``깨끗한 스킵'' 여부는, 응답의 \emph{공개 성분}(아래 참조)을 이용해
        공격자가 스스로 판정한다.
\end{description}

\paragraph{응답 벡터의 구조.} 레퍼런스에서 $c\svec_1$은 길이 $L=4$의 벡터로 계산되며, 그 0번째 성분은
$c$ 자체(비밀과 무관)이고 성분 1--3만이 $c\cdot\svec_1[0..2]$이다. 따라서 방출 응답 $z_1$은
1024계수($=4\times256$)이나 비밀은 768계수($=3\times256$)에만 담긴다. 성분 0은 $\pm c$로 공개되므로
부호와 \emph{깨끗한 스킵}(clean skip) 여부의 판정 기준이 된다.

\paragraph{실험 결과.} 위 절차를 T2 오류로 방출된 응답 $\zvec$에 적용한 결과, \textbf{$\svec_1$의 768개
계수($3\times256$)가 진짜 비밀키와 $\bmod\,q$에서 100\% 일치}함을 확인하였으며, 공개 블록 일치를 통한
자가 판정도 성립하였다. 식~\eqref{eq:t2recover}의 복원은 \texttt{HAETAE} 레퍼런스 정수 산술
(Montgomery/NTT)과 비트 단위로 일치한다(복구 파이프라인 self-test). 즉 가역 챌린지 조건 하에서 단 한
번의 T2 오류로, 다수의 서명 수집이나 통계적 분석 없이 비밀키 전체가 결정된다.

\subsection{오류 지점별 누설 방출률}
\label{subsec:repro}
각 오류 지점을 단일(one-shot) 오류로, 메시지를 바꿔 가며 결정론적으로 주입하여, 주입된 오류가 거부
샘플링을 통과해 누설 서명이 \emph{방출}되는 비율(누설 방출률)을 측정하였다(3개 배치 평균,
표~\ref{tab:repro}). 이 값은 오류 모델 하에서 지점별 \emph{논리적}
공격 난이도를 나타낸다.

\begin{table}[H]\centering
\caption{오류 지점별 누설 방출률(소프트웨어 오류 에뮬레이션; 거부 샘플링 통과 비율 = 지점별 공격 난이도).}
\label{tab:repro}
\small
\begin{tabular}{@{}lrrrl@{}}
\toprule
지점 & 누설 방출률 & 시도/누설 & 키 복원 & 특성 \\
\midrule
SEED, UNPACK        & 100\% & $\sim$1.0 & --- & 거부 무관 \\
SIGNBIT, LSB, T1    & $\sim$23\% & $\sim$4.3 & --- & 거부 의존 \\
T2                  & 83\% & $\sim$1.2 & \textbf{100\%} & 단일 서명 직접 \\
RB                  & 90\% & $\sim$1.1 & --- & 거부 판정 스킵 \\
\bottomrule
\end{tabular}
\end{table}

거부 의존 지점의 누설 방출률 $\sim$23\%는 선행 연구~\cite{lee2026haetae}가 보고한 거부 1회 통과 확률
$\sim$21\%와, 거부 무관 지점의 100\%와 각각 부합하여 본 실험 환경의 타당성을 교차 검증한다. T2는
83\%의 시도에서 누설이 방출되고, 방출된 모든 경우에 비밀키 768계수가 100\% 복원되어, 평균 1.2회의
오류만으로 완전한 단일 서명 키 복원이 가능함을 보인다(누설 방출률·복원은 소프트웨어 오류 에뮬레이션 하의 측정이다).

\section{제안 기법: IRV}
\label{sec:irv}

\subsection{설계 원리}
이중 연산이나 서명 후 검증과 같은 기존 대응은 모두 ``두 결과가 다르면 폐기한다'' 또는 ``검증에
실패하면 폐기한다''는 \emph{비교(분기) 연산}에 의존한다. 그러나 이러한 비교 연산 자체가
오류주입공격에 취약하며, 단일 오류로 해당 분기를 건너뛰면 방어가 무력화된다~\cite{yen2000}. 제안하는
IRV는 비교 연산에 의존하지 않고, 여러 무결성 검사의 잔차를 단일 누산기 $\delta$에 누적한 뒤
이를 \textbf{분기 없이} 출력 서명에 확산(감염)시킨다. $\delta=0$이면 서명은 그대로 유지되어
오탐이 없고, $\delta\neq0$이면 서명이 난수화되어 무효가 된다. ``오류가 발생하면 중단한다''는 조건
분기가 존재하지 않으므로, \emph{탐지 분기를 건너뛰는} 2차 오류가 성립하지 않는다(감염 연산 자체를
표적으로 하는 2차 오류에 대한 논의와 위협 모델의 범위는 \ref{sec:disc}절에서 다룬다).

\subsection{구성 요소}
제안 기법은 다음 무결성 검사(M1--M6, integrity-check Mechanism)를 응답 $\zvec$이 확정된 직후
\emph{서명 알고리즘 흐름 순서로} 수행하고, 모든 잔차를 $\delta$에 누적한다(알고리즘~\ref{alg:irv}). 여기서 각 검사의 \emph{잔차}(residual)란 재계산한 올바른 값과 실제 방출값의
불일치(XOR 또는 차)로, 오류가 없으면 $0$이고 오류가 있으면 $0$이 아니다. 이를 OR($\mathbin{|}$)로
누적하므로 $\delta=0$은 ``모든 검사 통과''와 정확히 동치이며(정상 실행에서 항상 성립), 이 경우
서명은 변경되지 않는다.
\begin{itemize}
  \item \textbf{M1} (비밀키 무결성): 비밀키 체크섬으로 서명 전 비밀키의 영속적 변조를 검출한다.
  \item \textbf{M2} (시드 정상성): 시드 버퍼가 0 또는 동일 패턴인지 검사하여 SEED를 차단한다.
  \item \textbf{M3} (부호$\cdot$난수 재유도): 채택된 시도의 $(\yvec,b)$를 재샘플링하여 대조함으로써
        SIGNBIT를 차단한다.
  \item \textbf{M4} (응답 재계산): $\zvec=\yvec+(-1)^b c\svec_1$을 재계산해 방출값과 대조하여(XOR 누적)
        T1과 T2를 \emph{비교 분기 없이} 차단한다.
  \item \textbf{M5} (서명경계 노름 재검사): $\lVert\zvec\rVert$가 \textbf{서명 경계 $B_1$과 bimodal 경계
        $B_0$}를 만족하는지 재검사하여 거부 우회(RB)를 차단한다.
  \item \textbf{M6} (표준 검증): 방출 서명을 검증식으로 재구성하여, 자기 일관적이지만 무효한
        서명(LSB, UNPACK)을 검출한다.
\end{itemize}
누적된 잔차 $\delta$는 마지막 \emph{감염 단계}에서 서명에 확산된다. $\delta\neq0$일 때만 전체 비트가 1이
되는 마스크 $m$를 분기 없이 생성하여\footnote{$m=-\big((\delta\mathbin{|}(-\delta))\gg(w{-}1)\big)$: $\delta\neq0$이면 $\delta$ 또는 $-\delta$의 최상위(부호) 비트가 1이라 $(w{-}1)$만큼 시프트하면 $1$, 음수화하면 전 비트가 $1$; $\delta=0$이면 $0$이다. 즉 조건 분기 없이 ``$\delta$가 0인가''를 전-$0$/전-$1$ 마스크로 변환한다. PRF에 서명별 값(메시지 또는 nonce)을 결합하는 것은 마스크 신선도를 위한 것으로, 동일 $\delta$가 여러 서명에서 반복되어도 마스크가 재사용되지 않게 한다.}, 서명에 $\mathrm{PRF}(\text{비밀 시드}\Vert\delta\Vert\text{메시지})\wedge m$를 XOR한다. 여기서 $\mathrm{PRF}$는 의사난수함수(pseudo-random function)이다.

\begin{algorithm}[H]
\caption{IRV 무결성 보호 (응답 $\zvec$ 확정 직후, 폭 $w$비트 누산)}
\label{alg:irv}
\begin{algorithmic}[1]
\State $\delta \gets 0$ \Comment{무부호 잔차 누산기}
\State $\delta \gets \delta \mathbin{|} \textsc{Chk}_{M1}(\svec\ \text{체크섬})$ \Comment{비밀키 무결성}
\State $\delta \gets \delta \mathbin{|} \textsc{Chk}_{M2}(\text{시드 정상성})$ \Comment{SEED}
\State $\delta \gets \delta \mathbin{|} \textsc{Chk}_{M3}(\yvec,b\ \text{재유도})$ \Comment{SIGNBIT}
\State $\delta \gets \delta \mathbin{|} \textsc{Chk}_{M4}\big(\zvec \overset{?}{=} \yvec+(-1)^b c\svec_1\big)$ \Comment{T1/T2}
\State $\delta \gets \delta \mathbin{|} \textsc{Chk}_{M5}\big(\lVert\zvec\rVert \le B_1 \ \wedge\ \text{bimodal } B_0\big)$ \Comment{RB}
\State $\delta \gets \delta \mathbin{|} \textsc{Chk}_{M6}(\text{표준 검증})$ \Comment{LSB/UNPACK}
\State $m \gets -\big((\delta \mathbin{|} (-\delta)) \gg (w-1)\big)$ \Comment{$\delta\!=\!0\Rightarrow m\!=\!0$, else 전비트 1 (무분기)}
\State $\sigma \gets \sigma \oplus \big(\mathrm{PRF}(\text{sk\_seed}\Vert\delta\Vert\text{msg}) \wedge m\big)$ \Comment{$\delta\!=\!0$이면 무변경}
\State \Return $\sigma$
\end{algorithmic}
\end{algorithm}

\paragraph{건전성과 잔존 위협.} (i)~\emph{오탐 없음}: 무결함 시 모든 검사 잔차가 0이라 $\delta=0$,
$m=0$이 되어 서명이 보존된다(\ref{subsec:coverage}절의 동일 다이제스트로 확인). (ii)~\emph{마스크
신선도}: 감염 마스크를 $\mathrm{PRF}(\text{비밀 시드}\Vert\delta\Vert\text{메시지})$로 유도하므로,
결정론 설정에서 동일 오류를 반복해도 메시지 결합으로 마스크가 신선하여 두 감염 출력의 XOR로 마스크가
상쇄되지 않는다. (iii)~\emph{우연 상쇄}: 서로 다른 오류가 동일 $\delta$를 유발할 확률은 누산 폭 $w$에
대해 $\approx 2^{-w}$로 무시할 수 있다. (iv)~\emph{잔존 2차 위협}: 본 기법은 \emph{탐지 분기} 스킵을
제거하나, 감염 적용(9행)이나 $\delta$ 누적(2--7행) 자체를 표적으로 하는 2차 오류는 여전히 단일 지점이며,
이에 대한 범위와 한계는 \ref{sec:disc}절에서 논한다.

\subsection{단순 서명 후 검증의 한계}
\label{subsec:b1b2}
직관적인 대응은 서명 후 표준 검증을 한 번 수행하는 것이다. 그러나 \texttt{HAETAE}의 검증 경계
$B_2$는 서명 경계 $B_1/B_0$보다 느슨하다($B_0,B_1<B_2$). 거부 판정 스킵으로 생성된 경계 초과
응답은 $B_1/B_0$는 위반하지만 $B_2$는 만족할 수 있어, 표준 검증을 그대로 통과한다. 따라서 서명 후
검증은 거부 우회(RB)를 검출하지 못한다. 반면 제안 기법의 M5는 검증 경계가 아니라 \emph{서명 경계}로
재검사하므로 이 사각지대를 차단한다. 즉 올바른 대응은 검증 경계가 아니라 서명 경계를 기준으로
재검사하여야 한다. 다만 서명 경로에는 $B_0$ bimodal 거부와 인코딩 상한(range coder)도 존재하여 경계
초과 응답의 일부는 인코딩 단계에서 재거부될 수 있으므로, RB 누설은 T2와 같은 즉각적 복원이 아니라
거부 샘플링 제거로 전사(transcript) 분포가 $\svec_1$에 의존하게 되는 \emph{통계적} 누설이다.

\section{평가}
\label{sec:eval}

본 장에서는 무방어 구현(\textsf{baseline})과 세 가지 대응 기법---전체 이중 연산(\textsf{double}),
선행 연구~\cite{lee2026haetae}의 대응을 재구현한 \textsf{leeha}, 그리고 제안 기법
\textsf{IRV}---을 동일한 조건에서 (i)~정확성, (ii)~커버리지, (iii)~2차 오류 내성,
(iv)~구현 비용의 네 가지 측면에서 비교한다.

\subsection{정확성}
\begin{itemize}
  \item 무결함 시 세 대응 기법 모두 \textsf{baseline}과 \textbf{동일한 서명(동일 다이제스트)}을 방출한다.
  \item \textsf{IRV}는 잔차 $\delta=0$이 보장되어 거짓 양성(오탐)이 발생하지 않는다.
\end{itemize}

\subsection{커버리지}
\label{subsec:coverage}
각 (오류, 변형) 쌍에 대해 오류를 주입하고, 방출된 서명의 다이제스트를 golden(무영향),
\textsc{leak}(\textsf{baseline}과 동일하여 미차단), blocked(무효화)의 세 가지로 분류한다
(표~\ref{tab:coverage}).

\begin{table}[H]\centering
\caption{오류 커버리지 및 2차 오류 내성. \texttt{blk}=blocked, \textsc{leak}=미차단, 우회=탐지 분기
스킵으로 무력화. $\dagger$ = 본 연구 도입 지점.}
\label{tab:coverage}
\small
\begin{tabular}{@{}lcccc@{}}
\toprule
오류 지점 & baseline & \textsf{double} & \textsf{leeha} & \textsf{IRV} \\
\midrule
SEED / SIGNBIT / UNPACK / LSB & \textsc{leak} & blk & blk & blk \\
T1 (CS)$^\dagger$ / T2 (ADDY)$^\dagger$ & \textsc{leak} & blk & blk & blk \\
RB (REJECT)$^\dagger$ & \textsc{leak} & blk & \textbf{\textsc{leak}} & blk \\
\midrule
차단 합계 / 7 & 0 & \textbf{7} & 6 & \textbf{7} \\
\midrule
2차 오류 (T2 + 탐지 분기 스킵) & \textsc{leak} & \textbf{\textsc{leak}(우회)} & \textbf{\textsc{leak}(우회)} & \textbf{차단} \\
\bottomrule
\end{tabular}
\end{table}

\textsf{double}과 \textsf{IRV}는 7개 지점을 모두 차단한 반면, \textsf{leeha}는 6개 지점을
차단하되 \textbf{거부 우회(RB)만 차단하지 못하였다}. RB는 본 논문의 위협 모델에 속하는 지점으로
선행 연구~\cite{lee2026haetae}의 범위 밖이다. 따라서 이 결과는 해당 연구의 오류로 해석해서는 안
되며, ``서명 후 검증 기법이 $B_1<B_2$로 인해 RB를 구조적으로 놓친다''는 의미로 읽어야 한다
(\ref{subsec:b1b2}절). 한편 \textsf{baseline}의 T2 \textsc{leak}이 실제 비밀키 누설임은
\ref{subsec:key-recovery}절에서 100\% 키 복원으로 입증한 바 있다.

\subsection{2차 오류 내성}
\label{subsec:2nd}
1차 오류(T2, 누설 생성)와 2차 오류(대응 기법의 탐지 분기 스킵)를 함께 주입하여, 탐지 분기를
무력화하는 공격에 대한 내성을 평가한다(표~\ref{tab:coverage} 마지막 행). 이는 표~\ref{tab:coverage}의
RB 커버리지(1차 오류 차단)와는 별개인 \emph{2차} 오류 위협에 대한 평가이다.

\textsf{double}의 결과 비교 분기와 \textsf{leeha}의 검증 분기는 모두 단일 스킵으로 우회되어, 누설
서명이 \textsf{baseline}과 동일하게 그대로 방출되었다. 반면 IRV는 응답 연산의 핵심 지점(T1/T2/RB)을
비교 분기 없이 감염형으로 처리하므로, 건너뛸 분기 자체가 존재하지 않아 차단을 유지하였다(상류 지점은
표준 검증으로 추가 보호된다). 다만 본 평가의 2차 오류는 \emph{탐지 분기 스킵} 모델에 한하며, 감염 연산
자체(마스크 적용$\cdot$잔차 누적)를 표적으로 하는 2차 오류는 \ref{sec:disc}절에서 별도로 논한다.

\subsection{비용}
\label{subsec:cost}
무결함 빌드의 코드/RAM(\texttt{arm-none-eabi-size})과 서명 1회 사이클(DWT)을 측정한다
(표~\ref{tab:cost}).

\begin{table}[H]\centering
\caption{구현 비용(무결함 빌드). 시간은 baseline 대비 상대값.}
\label{tab:cost}
\small
\begin{tabular}{@{}lccc@{}}
\toprule
변형 & 서명시간 & 코드(text) & 정적 RAM \\
\midrule
baseline       & $1.00\times$ & 24{,}556\,B & 31{,}168\,B \\
\textsf{double}& $2.00\times$ & $+0.7\%$    & $+0.6\%$ \\
\textsf{leeha} & $1.14\times$ & $+25.4\%$   & $+20.0\%$ \\
\textsf{IRV}   & $1.14\times$ & $+26.3\%$   & $+20.7\%$ \\
\bottomrule
\end{tabular}
\end{table}

이중 연산은 서명 시간이 2배로 증가하는 반면, \textsf{IRV}와 \textsf{leeha}는 각각 약 $1.14\times$에
그쳤다(표준 검증 1회와 상수 크기 검사를 포함). 코드 크기 증가 약 $25\%$는 주로 연결된 검증
루틴($\sim$6\,KB)에서 비롯된다. 즉 \textbf{서명 시간} 기준으로 IRV는 이중 연산의 $60\%$ 미만 비용으로
전면 커버리지와 거부 우회 차단, 2차 오류 내성을 동시에 달성한다. 다만 \textbf{코드/정적 RAM} 측면에서는
이중 연산(코드 $+0.7\%$, RAM $+0.6\%$)이 IRV($+26.3\%$, $+20.7\%$)보다 작으므로, IRV와 이중 연산은
(서명 시간$\cdot$2차 오류 내성) 대 (코드$\cdot$RAM)의 \emph{트레이드오프} 관계로 이해해야 한다.

\paragraph{\textsf{leeha} 비용에 관한 주석.}
공정한 비교를 위해 본 재구현의 \textsf{leeha}는 채택 이후 부분 이중 검사와 검증을 방출 서명당 한
번씩 수행하였으며, 이에 따라 $+14\%$의 비용이 발생하였다. 원 논문~\cite{lee2026haetae}은 더 가벼운
부분 이중 연산($+0.4\%$)과 서명 후 검증($+4.6\%$)을 보고하므로, 본 논문은 두 수치를 함께 제시한다.
다만 비용에 관한 핵심 결론은 이 세부 사항과 무관하다. 즉 IRV는 이중 연산의 절반 미만 비용으로
전면 커버리지와 거부 우회 차단, 2차 오류 내성을 모두 확보한다.

\subsection{종합}
\label{subsec:summary}
\begin{table}[H]\centering
\small
\begin{tabular}{@{}lcccl@{}}
\toprule
 & 커버리지 & RB & 2차오류 & 비용(시간 ; 코드/RAM) \\
\midrule
\textsf{double} & 7/7 & $\checkmark$ & 우회 & $2.0\times$ ; $+0.7\%/+0.6\%$ \\
\textsf{leeha}  & 6/7 & $\times$ & 우회 & $1.14\times$ ; $+25.4\%/+20.0\%$ \\
\textsf{IRV}    & \textbf{7/7} & \textbf{$\checkmark$} & \textbf{생존} & $1.14\times$ ; $+26.3\%/+20.7\%$ \\
\bottomrule
\end{tabular}
\end{table}

IRV는 (1)~거부 우회를 포함한 7개 지점 전면 차단, (2)~2차 오류 생존, (3)~이중 연산 절반 수준의
\emph{서명 시간}을 동시에 만족하는 유일한 변형이다. 두 대응 기법과의 비교를 나누어 보면 다음과 같다.
\begin{itemize}
  \item \textbf{vs \textsf{double}}: 커버리지는 7/7로 동일하나, IRV는 (i)~서명 시간이 절반($1.14\times$ 대
        $2.0\times$ --- IRV는 검증 1회, 이중 연산은 서명 전체를 2번 계산)이고, (ii)~2차 오류에 생존한다
        (이중 연산의 결과 비교 분기는 단일 스킵으로 우회된다). 대신 코드$\cdot$정적 RAM은 \textsf{double}이
        더 작아, 이 부분은 트레이드오프이다.
  \item \textbf{vs \textsf{leeha}}: 서명 시간은 비슷($\approx1.14\times$)하나, IRV는 (i)~거부 우회(RB)를
        차단하고 (ii)~2차 오류에 생존한다는 점에서 앞선다 --- \textsf{leeha}는 둘 다 놓친다(원 논문의
        경량 구성은 \ref{subsec:cost}절 주석 참조).
\end{itemize}

\subsection{논의 및 한계}
\label{sec:disc}

본 절에서는 제안 기법과 실험의 한계를 명확히 하고, 이를 둘러싼 해석상의 주의점을 정리한다.

\begin{itemize}
  \item \textbf{선행 기법의 차용.} IRV의 표준 검증(M6), 부분 이중 검사(M3), 정상성 검사(M2)는
        선행 연구~\cite{lee2026haetae}와 탐지 원리를 공유한다. 본 논문의 기여는 새로운 탐지 항목을
        제시한 데 있지 않고, 이들을 \emph{비교 연산 없는 감염형으로 통합}하여 2차 오류 내성을 부여하고,
        서명 경계 재검사로 거부 우회를 차단하며, 단일 서명 T2 공격을 오류 모델 하에서 실증한 데 있다.
  \item \textbf{거부 우회의 범위.} RB의 \textsc{leak}은 ``서명 경계를 위반한 응답이 검증을
        통과한다''는 의미이며, T2와 같은 즉각적 키 복원이 아니라 통계적 누설에 해당한다.
  \item \textbf{\textsf{leeha} 비용.} \ref{subsec:cost}절의 수치는 공정 비교를 위한 재구현 결과이며,
        원 논문의 보고치와 함께 제시하였다.
  \item \textbf{오류 모델의 현실성.} 본 논문이 가정한 명령어 스킵은 32비트 마이크로컨트롤러에서
        물리적으로 실증된 표준 오류 모델이며~\cite{moro2013}, 소프트웨어 라인 주입은 이 오류의
        \emph{논리적} 효과를 결정론적으로 재현하여 대응 기법을 공정하게 비교한다. 본 논문의 공격 및
        대응 평가는 이 에뮬레이션된 단일 오류 모델을 기준으로 하며, 물리적 오류 주입을 통한 실현은
        향후 과제로 남긴다.
  \item \textbf{2차 오류 위협 모델.} 본 논문은 출력을 게이팅하는 \emph{조건 분기} 1회 스킵을
        모델링하며, IRV의 생존은 응답 연산을 비교 분기 없이 처리한 데서 비롯된다. 이 모델은
        detect-and-abort의 구조적 약점(뒤집을 불리언 분기의 존재)을 정확히 겨냥한 것으로, IRV는
        그러한 분기를 제거한다. 다만 동일 능력의 공격자는 감염 적용 단계(마스크 XOR)나 잔차 누적
        자체를 표적으로 하는 2차 오류를 시도할 수 있으며, 이 경우 감염형 대응도 무력화될 수 있다.
        즉 본 논문의 2차 오류 내성 주장은 \emph{탐지 분기 스킵 모델에 한정}되며, 감염 연산을
        데이터 흐름에 직조하여 이러한 잔존 위협까지 견디게 하는 강건 구현은 향후 과제다.
  \item \textbf{대리 주입과 결정론 설정.} 소프트웨어 라인 주입과 결정론 난수 발생기는 대응 기법을
        공정하게 비교하기 위한 설정이며, 복구 절차는 소프트웨어 클린 스킵 누설에 대한 self-test로 100\% 검증하였다.
  \item \textbf{절대 수치의 키 의존성.} \textsf{baseline}의 60.86\,M 사이클은 해당 키의 거부 반복
        횟수에 의존하므로, 본 논문은 \emph{상대} 비용을 기준으로 결론을 제시한다.
  \item \textbf{일반화.} 제안 원리는 HAETAE 120/180/260은 물론, 동일한 응답 구조를 갖는 Fiat--Shamir
        with aborts 계열(\texttt{Dilithium}~\cite{dilithium} 등)로 확장할 수 있다.
\end{itemize}

\section{관련 연구}
\label{sec:related}

\textbf{격자 서명에 대한 오류주입공격.} Espitau 등~\cite{espitau2016}은 거부 샘플링 루프를 조기
종료시키는 loop-abort 공격을 제안하였고, Groot Bruinderink와 Pessl~\cite{bruinderink2018}은 결정론
격자 서명에 대한 클럭 글리치 차분 공격을 보고하였다. Ravi 등~\cite{ravi2022survey}은 Kyber와
Dilithium에 대한 부채널 및 오류 공격을 종합하였으며, Dilithium에 대해서는 정정 오류와 명령어 스킵을
이용한 공격이 실험적으로 입증되었으며~\cite{krahmer2024,elghamrawy2023}, 특히 Ravi
등~\cite{ravi2019determinism}은 결정론 격자 서명에서 덧셈 스킵(skip-addition) 오류로 비밀키의
일부를 복원하였다. 또한 NTT 트위들 상수에 대한 오류 분석~\cite{ravi2023ntt}과 hedged ML-DSA에 대한
단일 서명 오류공격~\cite{jendral2024}이 보고되었다. 본 논문은 \texttt{HAETAE}에서 $+\yvec$
스킵(T2)이 단일 서명만으로 비밀키 \emph{전체}를 직접 복원시킴을 새롭게 제시한다(표~\ref{tab:attack-compare}).
기법의 핵심 착상($\yvec$ 제거 후 가역 챌린지로 나눔)은 loop-abort~\cite{espitau2016}·skip-addition
~\cite{ravi2019determinism} 계열과 계보를 공유하나, 본 논문의 차별점은 $+\yvec$ 전체 스킵으로
\emph{단일} 서명에서 $\svec_1$ \emph{전체}를 직접 복원한다는 점이다.

\begin{table}[H]\centering
\caption{격자 서명 오류주입공격 비교(오류 서명 수 = 키 복원에 필요한 결함 서명 개수).}
\label{tab:attack-compare}
\small
\begin{tabular}{@{}lllc@{}}
\toprule
연구 & 대상 & 오류 효과 & 오류 서명 수 \\
\midrule
Espitau 등~\cite{espitau2016}             & BLISS/GLP          & 거부 루프 조기 종료          & 다중 \\
Groot Bruinderink--Pessl~\cite{bruinderink2018} & 결정론 서명   & 클럭 글리치 차분             & 다중 \\
Ravi 등~\cite{ravi2019determinism}        & \texttt{Dilithium} & 덧셈 스킵(\emph{부분} 키)    & 다중 \\
Lee 등~\cite{lee2026haetae}               & \texttt{HAETAE}    & 부호 비트 등 4지점(차분)     & $\ge 2$ \\
\textbf{본 논문 (T2)}                     & \texttt{HAETAE}    & $+\yvec$ 스킵 $\to \zvec{=}\pm c\svec_1$ & \textbf{단일} \\
\bottomrule
\end{tabular}
\end{table}

\textbf{양자내성암호 대응 기법.} 대표적인 대응으로 이중 연산(높은 비용), 서명 후 검증(경량이나
거부 우회를 놓침), 셔플링(비결정적)이 있다. 프로세서 수준 및 소프트웨어 수준의 오류 대응은 폭넓게
연구되어 왔으나~\cite{boulifa2025,barenghi2010}, 대부분 비용과 안전성 사이의 절충에 직면한다.
제안 기법은 서명 경계 재검사로 서명 후 검증의 사각지대를 메우고, 비교 연산 없는 감염형 처리로
셔플링의 비결정성을 피하며, 이중 연산보다 낮은 비용으로 동작한다. 특히 셔플링은 Ulitzsch
등~\cite{ulitzsch2023}이 정수 선형 계획법(ILP)으로 무력화함을 보인 바 있어, 오류에 대한 근본적
대응으로는 한계가 있다.

\textbf{선행 HAETAE 오류 연구.} Lee 등~\cite{lee2026haetae}은 \texttt{HAETAE}의 네 가지 공격
지점(LSB, 공개 행렬 언패킹, 부호 비트, 샘플링 시드)과 네 가지 대응(서명 후 검증, 거부 루프 내
이동, 부분 이중 연산, 정상성 검사)을 제시하였다. 본 논문은 이를 직접적인 비교 대상으로 삼되, 겹치는
탐지 항목은 명시적으로 인용한다.

\textbf{감염형 대응의 계보.} 본 논문은 Fermat 소정리에 기반하여 비교 연산 없이 RSA 오류를 검출한
Seo 등~\cite{seo2013}의 대응 기법을 격자 Fiat--Shamir 서명으로 확장한 것이다. RSA 오류 대응의
형식 검증~\cite{christofi2013}과 격자 서명의 오류 민감도 분석~\cite{bindel2016}도 본 연구와 관련된다.

\section{결론}
\label{sec:concl}

본 논문에서는 한국형 양자내성암호 서명 후보인 \texttt{HAETAE}의 응답 연산
$\zvec=\yvec+(-1)^b c\svec_1$에 대한 단일 오류주입공격을 분석하고, $+\yvec$ 덧셈을 건너뛰는
경우(T2) 단 한 번의 오류만으로 비밀키 $\svec_1$을 직접 복원할 수 있음을 STM32F405 전체 서명
환경에서 \textbf{오류 모델 하에} 실증하였다. 또한 이를
효율적으로 방어하기 위해, 연산 재계산 검증과 서명 경계 노름
재검사, 비밀키 무결성 검사를 \textbf{비교 연산에 의존하지 않는 감염형 마스킹}으로 통합한 경량 대응
기법 \textbf{IRV}를 제안하였다. 제안 기법은 거부 우회를 포함한 7개 오류 지점을 모두
차단하고(탐지 분기 스킵 모델 하) 2차 오류에 내성을 가지며, 이중 연산의 절반 미만 \emph{서명 시간}으로
동작함을 실험적으로 입증하였다. 이로써 IRV는 서명 시간과 안전성(거부 우회 차단$\cdot$2차 오류 내성)
측면에서 이중 연산 및 서명 후 검증보다 유리하다(코드$\cdot$정적 RAM은 이중 연산이 더 작다). 향후
연구로는 T1 차분(서명 2개) 공격과 동일 구조를 갖는 다른 격자 서명으로의 확장을 계획한다.


\newpage
\appendix
\section{IRV 잔차 계산과 감염의 수치 예제}
\label{app:irv-example}

본 부록은 알고리즘~\ref{alg:irv}의 잔차 누산기 $\delta$와 각 검사 $\textsc{Chk}_{Mi}$가 구체적으로
어떻게 계산되는지, 그리고 $\delta$가 어떻게 \emph{분기 없이} 서명에 감염되는지를 작은 수치 예로 보인다.

\subsection{잔차와 $\textsc{Chk}$의 정의}
각 검사 $\textsc{Chk}_{Mi}$는 (i)~올바른 값을 \emph{재계산}하고, (ii)~서명에 실제로 방출된 값과
\emph{XOR}(또는 차)을 취한다. 두 값이 같으면 잔차는 $0$, 다르면 $0$이 아니다. 6개 검사의 잔차를 비트
OR로 누적하므로
\[
  \delta \;=\; \textsc{Chk}_{M1} \mathbin{|} \textsc{Chk}_{M2} \mathbin{|} \cdots \mathbin{|} \textsc{Chk}_{M6}
\]
가 되며, $\delta=0$은 ``모든 검사 통과''와 정확히 동치이다. 이어 폭 $w$비트 마스크
\[
  m \;=\; -\big((\delta \mathbin{|} (-\delta)) \gg (w-1)\big)
\]
는 $\delta=0$이면 $m=0$(전-$0$), $\delta\neq0$이면 $m$이 전-$1$이 되어, 서명에
$\mathrm{PRF}(\text{sk\_seed}\Vert\delta\Vert\text{msg}) \wedge m$을 XOR한다. 즉 무결함이면 서명이 그대로
보존되고, 결함이면 서명 전체가 난수화되어 무효가 된다 --- \emph{조건 분기 없이}.

\subsection{응답 재계산 검사 $\textsc{Chk}_{M4}$의 수치 예}
설명을 위해 응답의 한 계수만 본다. 어떤 계수에서 $y_1=100$, $c\svec_1=7$, 부호 $(-1)^b=+1$이라 하면
올바른 응답은 $z_1 = 100+7 = 107$이다. $\textsc{Chk}_{M4}$는 방출된 $z_1$과 재계산값 $100+7=107$을
XOR한다.

\paragraph{무결함.} 방출값 $107$, 재계산값 $107$ $\Rightarrow$ $107 \oplus 107 = 0$. 따라서 이 계수의
잔차는 $0$이고, 다른 검사도 모두 $0$이면 $\delta=0$, $m=0$이 되어 서명은 \textbf{변경되지 않는다}.

\paragraph{T2 결함($+\yvec$ 스킵).} 이때 방출값은 $y$가 빠진 $z_1 = 7$이나 재계산값은 여전히
$100+7=107$이다. $\Rightarrow$ $7 \oplus 107 = 108 \neq 0$. 따라서 $\delta \mathrel{|}= 108$로
$\delta\neq0$이 되고,
\[
  m = -\big((108 \mathbin{|} (-108)) \gg (w-1)\big) = \underbrace{11\cdots1}_{w\text{비트}}
\]
이므로 서명은 $\mathrm{PRF}(\cdots)\wedge m = \mathrm{PRF}(\cdots)$와 XOR되어 \textbf{난수화(무효)}된다.
즉 T2로 만든 누설 서명이 그대로 방출되지 못한다.

\subsection{다른 검사의 잔차}
같은 방식으로 나머지 검사도 잔차를 낸다: $\textsc{Chk}_{M1}$은 비밀키 체크섬을 재계산해 저장값과 XOR,
$\textsc{Chk}_{M2}\cdot\textsc{Chk}_{M3}$은 시드/재유도한 $(\yvec,b)$의 불일치 비트,
$\textsc{Chk}_{M5}$는 $\lVert\zvec\rVert^2$이 경계 $B_1^2 L_n^2$을 넘는지의 부호 비트(위반 시 $1$),
$\textsc{Chk}_{M6}$은 표준 검증 실패 비트를 각각 잔차로 낸다. 어느 하나라도 $0$이 아니면 $\delta\neq0$이
되어 위 A.2와 동일하게 서명이 무효화된다. 그러므로 IRV는 7개 오류 지점 중 무엇이 발생하든 \emph{비교
분기 없이} 누설 서명의 방출을 막는다.


\newpage
\bibliographystyle{plain}
\bibliography{references}

\end{document}
